\section{Uniform analytic estimates for the linear block map}
\label{app:analytic-block}

This section supplies the analytic estimates used in the linear block
transformation.  In particular, every strip width and every constant below
is chosen independently of the blocking depth and before the sufficiently
large integer $L$ is chosen.  There are only finitely many derivative and
decay orders in an application of these estimates.

\subsection{Domains, sampling, and the free precision}

We use integer site coordinates for Fourier transforms.  Changing to
block-centered coordinates conjugates the formulas by deterministic
translation phases; in coarse momentum these phases and their fixed
derivatives are uniformly bounded.  The integer-coordinate convention
makes the gluing of aliases explicit.  For $\delta>0$, put
\[
 \Omega_\delta=\{p\in\mathbb C^2:
       |\mathop{\rm Re}p_\mu|<\pi+\delta,\quad
       |\mathop{\rm Im}p_\mu|<\delta\}.
\]
An alias field is a family on these overlapping cubes, with the alias
index relabeled when $p$ crosses a face of the real momentum cube.
Equivalently it is a function of $k$ on the fine momentum torus, with
$p=Lk$ reduced modulo $2\pi$.  At a centered representative $l$ write
$\langle l\rangle=1+|l|_1$.  The weights of two representatives used
on an overlap are comparable by an absolute constant.  All estimates
on overlapping cubes therefore define estimates of the same field.
An analytic star means $f^*(z)=\overline{f(\bar z)}$ for real-kernel
matrix symbols, with transposition included for matrices; for the
scalar Fourier transform of real weights it is $b^*(z)=b(-z)$.

\begin{lemma}[Uniform shells and strips]\label{app:free-strip}
For the sampled smooth block weights in Section~\ref{free:section},
there are $\delta,c,C>0$ and $L_0$ with the following properties for
$L\ge L_0$ and $h\ge0$.  The precisions $K_h$ and an elliptic edge lift
$B_h$ are analytic on the complex $\delta$-neighborhood of the momentum
torus.  On the real torus,
\[
 cD\le K_h\le CD,\qquad cI\le B_h\le CI,
 \qquad K_h=d^*B_hd,\quad B_h(0)=I.
\]
Their fixed derivatives are bounded, $K_h-D$ has a zero of order four
at zero, and their history differences, including the inverse edge
lifts, are bounded by $CL^{-h}$ when compared with any depth
$\bar h\ge h$.  The same assertions hold on a common smaller strip.
\end{lemma}

\begin{proof}
Let $s_L=\sum_xb_x^2$.  Smoothness and normalization of the sampled
bump give $s_L\le A/L^2$ for a fixed $A$.  Discrete summation by parts
$N$ times, with no boundary terms because the bump is supported in the
block interior, gives
\begin{equation}
 |\partial_p^\alpha b((p+2\pi l)/L)|
       \le C_{N,\alpha}\langle l\rangle^{-N}
       \quad(p\in\Omega_{\delta_0}).                 \label{app:analytic-1}
\end{equation}
Here $\delta_0>0$ can initially be any sufficiently small fixed
constant.  Indeed the sum of absolute $N$th differences of the
normalized weights is $O(L^{-N})$, whereas on a nonzero alias at
least one difference multiplier has size $c\langle l\rangle/L$.
Complex tilting costs at most $\exp(C\delta_0)$, since the block has
diameter $O(L)$; derivatives in $p$ insert bounded rescaled coordinates.
The principal alias is bounded directly.  This also proves \eqref{app:analytic-1} in
the centered convention.

Set $m=L^h$, $W_h(\xi)=\prod_{j=1}^h b(\xi/L^j)$ and
$D_m(\xi)=m^2D(\xi/m)$.  For centered $l\ne0$ modulo $m$,
\begin{equation}
 |D_m(p+2\pi l)|\ge c\langle l\rangle^2
                    \quad(p\in\Omega_{\delta_0}),       \label{app:analytic-2}
\end{equation}
after reducing $\delta_0$.  To check the complex bound directly,
use $D_m(z)=4m^2\sum_\mu\sin^2(z_\mu/(2m))$.
Its real part is at least $c|\mathop{\rm Re}z|^2-C|\mathop{\rm Im}z|^2$
on this cube.  A nonzero alias has
$|\mathop{\rm Re}(p+2\pi l)|\ge c\langle l\rangle$.
This proves \eqref{app:analytic-2}, including aliases at the faces of the cube.

Here is a precise shell estimate, including its complex tilt.  The
$r$-fold block weights have squared sum $s_L^r$: the base-$L$
representation of each position in an $L^r$ block is unique.  Their
positions divided by $L^r$ remain in a fixed bounded square.  Finite
Fourier orthogonality, first at tilt $+\mathop{\rm Im}p$ and then at
tilt $-\mathop{\rm Im}p$, and Cauchy--Schwarz give
\begin{equation}
 \sum_{l\bmod L^r}
 |W_r(p+2\pi l)W_r(p+2\pi l)^*|\le C A^r.             \label{app:analytic-3}
\end{equation}
For a shell contained in $|l|_\infty\le L^r/2$, the factors with
$j>r$ have absolute product at most a fixed constant: their real
arguments cost at most one and their complex tilts cost
$\exp(C\delta_0\sum_{j>r}L^{1-j})$.  Thus \eqref{app:analytic-3} bounds the
corresponding shell of $W_h$ too.  Split the nonprincipal aliases
into the first shell $0<|l|_\infty\le L/2$ and the shells
$L^{r-1}/2<|l|_\infty\le L^r/2$, $2\le r\le h$.
Equations \eqref{app:analytic-2}--\eqref{app:analytic-3} imply
\begin{equation}
 \sum_{l\ne0}
 \left|\frac{W_h(p+2\pi l)W_h(p+2\pi l)^*}
 {D_m(p+2\pi l)}\right|
 \le C\sum_{r=1}^h\frac{A^r}{L^{2(r-1)}}\le C,        \label{app:analytic-4}
\end{equation}
uniformly once $L^2\ge2A$.  The noise sum
$a^{-1}\sum_{r<h}s_L^r$ is also bounded.  Consequently the function
$R_h$ in \eqref{free:explicit} is bounded on a fixed complex domain.
Cauchy's estimate on nested fixed domains bounds every required fixed
derivative; no real-momentum estimate is being used in place of a
complex estimate.

For real $p$ in the cube, the centered one-dimensional bump has phase
variation less than $\pi/2$ in its first principal factor.  The
product structure gives $|b(p/L)|\ge c$.
For $j\ge2$,
$1-b(p/L^j)b(p/L^j)^*=O(|p|^2L^{-2(j-1)})$.
It follows that $W_h(p)W_h(p)^*\ge c$ uniformly in $h$.
The denominator
\[
 F_h(p)=W_h(p)W_h(p)^*+D_m(p)R_h(p)
\]
is therefore at least $c$ on the real cube.  It and its first
derivatives have uniform bounds on the larger complex cube.  A
fixed reduction of the imaginary width gives $|F_h|\ge c/2$.
The formula $K_h=D_m/F_h$ now proves analyticity and boundedness.
The defining covariance sum proves that expressions from overlapping
alias cubes agree.  Real ellipticity follows from \eqref{app:analytic-4}, and the
expansions $F_h=1+O(|p|^2)$,
$D_m=|p|^2+O(|p|^4/m^2)$ give $K_h-D=O(|p|^4)$.

For completeness, history differences require no division by a bump
factor.  On common aliases $|l|_\infty\le L^h/2$ put
$\xi=p+2\pi l$.  Taylor's theorem for the extra product, with the
linear term canceled in each $bb^*$, gives
\[
 |W_{\bar h}W_{\bar h}^*-W_hW_h^*|
       \le C|\xi|^2L^{-2h}|W_hW_h^*|.
\]
On the same aliases,
$|D_{L^{\bar h}}^{-1}-D_{L^h}^{-1}|\le CL^{-2h}$;
this follows from
$|D_{L^{\bar h}}-D_{L^h}|\le C|\xi|^4L^{-2h}$ and \eqref{app:analytic-2}.
Thus the common nonprincipal sum changes by at most
$CA^hL^{-2h}$.  The extra shells have total at most
$CA^{h+1}L^{-2h}$ by \eqref{app:analytic-4}, and the noise tail has the same bound.
After increasing $L_0$ these quantities are at most $CL^{-h}$.
The principal product and numerator have their direct Taylor bounds.
The nonvanishing denominator then gives the history estimate for
$K_h$, and Cauchy's estimate gives its fixed derivatives.

We spell out why the edge lift does not introduce a scale-dependent
strip.  If $f$ is periodic analytic, the quotient
\[
 \frac{f(p_1,p_2)-f(0,p_2)}{e^{ip_1}-1}
\]
is analytic and bounded on a smaller fixed strip, with norm bounded
by a fixed multiple of the norm of $f$ on the larger strip.
Near $p_1=0$ this is Taylor's integral formula; away from zero the
denominator is bounded below.  Repeat in the two coordinates.
If $f$ vanishes through degree three at zero this writes it as a
sum of the five monomials $d_1^j d_2^{4-j}$ with bounded analytic
coefficients.  Since $d_\mu^*=-e^{-ip_\mu}d_\mu$, this gives
$K_h=d^*B_h^{\rm raw}d$, with
$B_h^{\rm raw}=I+O(|p|^2)$ and uniformly bounded coefficients.
Taking its Hermitian, reality, and lattice-symmetry averages preserves
the identity.  Put $c(p)=(-d_2,d_1)$.
Adding $\lambda c^*c$ leaves the identity unchanged.  Near zero
$B_h^{\rm raw}$ is uniformly positive.  Away from zero its
longitudinal quadratic form is $K_h/D\ge c$, its mixed entries are
bounded, and the transverse addition is $\lambda D$.
The elementary two-by-two Schur complement therefore makes the
matrix uniformly positive for one fixed $\lambda$.
Uniform first-derivative bounds make this matrix and its inverse
analytic on a smaller fixed strip.  All operations before this
common addition are bounded linear operations; the addition cancels
in differences.  The inverse difference identity supplies the claimed
history bound.  Depth zero can be kept equal to $I$; its discrepancy
bound is only the bounded $h=0$ bound.  This proves the lemma.
\end{proof}

\subsection{The two divisions in the edge intertwiner}

The following elementary analytic fact specifies the divisions needed
in the construction of $T$.  It is useful to keep the alias scale
$L^{-1}$ in its statement.

\begin{lemma}[Alias syzygy and coarse division]\label{app:alias-division}
Suppose an analytic alias vector $R_l(p)$ satisfies
$d(k_l)^*R_l(p)=0$, has the bounds
\[
 |R_l(p)|\le C_sL^{-1}\langle l\rangle^{-s}
\]
on a fixed larger domain for every required fixed $s$, and has a zero
through first order at the principal origin.  Then
$R_l=c(k_l)^*\gamma_l$ on a smaller fixed domain, with
$|\gamma_l|\le C_s\langle l\rangle^{-s}$ and
$\gamma_0(0)=0$.  If $R_l(0)=0$ on every nonprincipal alias, then
$\gamma_l(0)=0$ for all $l$.

If, in addition, $\gamma_l(0)=0$ for every alias, there is an
analytic row $H_l$ satisfying
$H_ld'(p)=\gamma_l(p)$ and
$|H_l|\le C_s\langle l\rangle^{-s}$.  Both constructions are linear,
preserve alias gluing, and have constants independent of $L$.
\end{lemma}

\begin{proof}
On a nonprincipal alias one of $|d_\mu(k_l)^*|$ is at least
$c\langle l\rangle/L$ throughout the domain.  Divide the
corresponding component of $R$ by that component of $c^*$.
Where both choices are possible they agree by $d^*R=0$.
This gives the asserted estimate with one extra power of decay,
which can be discarded.  On the principal alias away from a small
fixed polydisc the same argument applies with denominator $c/L$.
Inside that polydisc the identity implies
$R_2(0,p_2)=0$ and $R_1(p_1,0)=0$.
Hence $R_2/d_1^*=-R_1/d_2^*$ extends analytically, by the coordinate
division proved above.  Since $d_\mu(p/L)^*$ is $p_\mu/L$ times
an analytic unit, its norm is bounded by $CL\|R\|$, with a fixed
Cauchy constant.  The first-order vanishing of $R$ makes the
quotient zero at the origin.  At a nonprincipal zero-momentum alias
the nonzero divisor gives the last assertion directly.  Uniqueness
away from the common zero proves agreement on every overlap.

We give the second division explicitly to avoid summing $L$ bounded
interpolation terms.  Fix the largest desired decay order $s$ and a
single even integer $J>s+2$, increasing it further if required by
the finite list of difference orders.  On the fine torus set
\[
 (\Pi\gamma)(k)=\sum_{n\bmod L}
 \left(\frac{\sin(L(k_1-2\pi n/L)/2)}
 {L\sin((k_1-2\pi n/L)/2)}\right)^J
                 \gamma(2\pi n/L,k_2).
\]
The power is periodic because $J$ is even.  It is one at its own
sampling point and zero at all the others.  On alias $l$ its
absolute value is at most
$C_J(1+\operatorname{dist}_{\rm cyc}(l_1,n))^{-J}$
on a fixed complex cube: the numerator is bounded there and the
denominator is bounded below by a constant times the cyclic
distance except at the removable own sample, where the quotient
is bounded directly.  The elementary triangle inequality gives
\begin{equation}
 \langle l\rangle^s
 \sum_{n\bmod L}
 \frac{\langle(n,l_2)\rangle^{-s}}
 {(1+\operatorname{dist}_{\rm cyc}(l_1,n))^J}
 \le C_s\sum_{j\in\mathbb Z}(1+|j|)^{s-J}\le C_{s,J}.
                                                               \label{app:analytic-5}
\end{equation}
Thus $\Pi$ is bounded in the same weighted alias norm, independently
of $L$.  Its fixed coarse derivatives are bounded by Cauchy's
estimate on nested fixed cubes.  Now set
\[
 H_1=\frac{\gamma-\Pi\gamma}{e^{ip_1}-1},\qquad
 H_2=\frac{\Pi\gamma}{e^{ip_2}-1}.
\]
The first numerator vanishes on $p_1=0$ by interpolation.  On
$p_2=0$, all samples in the second numerator vanish by
$\gamma_n(0)=0$.  Coordinate division in $p$, rather than in $k$,
therefore bounds both quotients with no power of $L$.
The resulting identity is $H_1d'_1+H_2d'_2=\gamma$.
Only a fixed finite number of domain reductions has been used.
\end{proof}

We verify the hypotheses for the actual construction.  On a
nonprincipal alias, Lemma~\ref{app:free-strip} and its explicit
formula imply
$|K_h(k_l)^{-1}|\le CL^2\langle l\rangle^{-2}$
on the coarse complex cube.  The bump estimate \eqref{app:analytic-1} and
\eqref{free:P} give arbitrary fixed alias decay for $P_l$.
On the principal alias the exact identity
\[
 P_0=b(k_0)^{-1}\left[1-
 \left(a^{-1}+L^{-2}\sum_{l\ne0}b(k_l)b(k_l)^*K_h(k_l)^{-1}
 \right)K_{h+1}(p)\right]
\]
removes the pole.  The principal bump is nonzero on the fixed
smaller domain by the preceding principal-factor estimate.
Consequently $d_lP_l$ has bound
$C_sL^{-1}\langle l\rangle^{-s}$.
The finite geometric sum defining $\omega_{\mu,l}$ is bounded
by $CL$ on that domain; hence $T_{0,l}$ has the same bound.
The identities $KP=Q^*K'$ and $\omega_\mu d_\mu=d'_\mu$
give $d_l^*R_l=0$ exactly.  At $p=0$ both terms in $R_l$ vanish
on nonprincipal aliases.  On the principal alias their linear
terms are both $ip/L$, so $R_0$ vanishes through first order.
Lemma~\ref{app:alias-division} therefore gives the stated $\gamma,H$.
In
\[
 T_l=T_{0,l}+c_l^*H_l(B')^{-1}
\]
the extra factor $c_l^*$ has size $C\langle l\rangle/L$;
using one further decay order proves
$|T_l|\le C_sL^{-1}\langle l\rangle^{-s}$.
The construction is linear in $R$ and uses uniformly bounded analytic
inverses, so its history differences have the same $L^{-h}$ factor.
Choose $J$ once after listing all required decay and difference orders.

Finally, the remaining factorizations in the positive completion
introduce no further scale loss.  A coarse analytic row $v$ with
$vd'=0$ has the form $v=\eta c'$ by the same coordinate-plane
division, with bounded $\eta$.  A coarse analytic matrix $A$ with
$Ad'=0$ and $d'^*A=0$ consequently has the form
$A=c'^*\eta c'$, by applying that division successively to its
columns and rows.  The estimates use only two fixed Cauchy constants.
Apply this to $B'UB'$ and to $c_lB_l^{-1}T_lB'$.
The latter row before division is
$O(L^{-2}\langle l\rangle^{-s})$.
Its squared fine norm includes the factor $L^2\sum_l$, and hence
is $O(L^{-2})$ after division.  These are precisely the bounds used
to choose the curl correction with its constant fixed before $L$.
Fourier inversion on any smaller retained strip, and the factors
$d(k_l)$ for fine differences, now give the asserted exponential
kernel moments.  All reductions of width in this section are made
once for the construction; they are not repeated as the depth grows.
